\emph{পূর্ণতা উপপাদ্য} \emph{বিশুদ্ধতা উপপাদ্যের} বিপরীত।
এক রূপে এটি বলে যে $\Gamma \Entails !A$ হলে
$\Gamma \Proves !A$। অন্য রূপে বলে, $\Gamma$ সঙ্গত হলে সেটি
পরিতৃপ্তিযোগ্য। আমরা দ্বিতীয় রূপটি প্রমাণ করেছি (এবং দ্বিতীয়টি
থেকে প্রথমটি প্রতিপাদন করেছি)। প্রমাণটি জটিল এবং এতে কয়েকটি
ধাপ লাগে। একটি সঙ্গত সেট~$\Gamma$ দিয়ে শুরু করি। প্রথমে
অসীমসংখ্যক নতুন ধ্রুবক সংকেত~$c_i$ এবং
$\lexists[x][!A(x)] \lif !A(c)$ আকারের সূত্র যোগ করি। এতে
বর্ধিত ভাষার যে প্রতিটি সূত্রে একটি মুক্ত চলরাশি আছে, সেই
সূত্র~$!A(x)$-এর সঙ্গে
নতুন ধ্রুবকগুলির একটি জুড়ে দেওয়া হয়। এর ফলে~$\Gamma$-কে
ধারণ করে এমন একটি \emph{সম্পৃক্ত} সঙ্গত বাক্যসমষ্টি পাওয়া
যায়। সেটটি তখনও সঙ্গত। এবার সেই সেটটিকে প্রসারিত করে একটি
\emph{পূর্ণ সঙ্গত সেট} তৈরি করি। পূর্ণ সঙ্গত সেটের একটি সুবিধাজনক
ধর্ম হলো, যেকোনো বাক্য~$!A$-এর জন্য হয় $!A$ নয়তো $\lnot !A$
সেটটিতে থাকে (দুটিই কখনো থাকে না)। যেহেতু একটি সম্পৃক্ত সেট
থেকে শুরু করেছিলাম, এখন আমাদের কাছে বাক্যের একটি সম্পৃক্ত,
পূর্ণ, সঙ্গত সেট~$\Gamma^*$ আছে, যাতে~$\Gamma$ অন্তর্ভুক্ত।
এবার এই সেট থেকে একটি গঠন~$\Struct{M}$ সংজ্ঞায়িত করা যায়,
যাতে $\Sat{M(\Gamma^*)}{!A}$ হয় ঠিক তখনই, যখন
$!A \in \Gamma^*$। বিশেষত, $\Sat{M(\Gamma^*)}{\Gamma}$, অর্থাৎ
$\Gamma$ পরিতৃপ্তিযোগ্য। $=$ উপস্থিত থাকলে নির্মাণটি কিছুটা
বেশি জটিল।

পূর্ণতা উপপাদ্য থেকে দুটি গুরুত্বপূর্ণ অনুসিদ্ধান্ত পাওয়া যায়।
\emph{সংহতি উপপাদ্য} বলে যে $\Gamma \Entails !A$ হয় ঠিক তখনই,
যখন $\Gamma_0 \Entails !A$ হয় কোনো সসীম
$\Gamma_0 \subseteq \Gamma$-এর জন্য। এর একটি সমতুল্য বক্তব্য হলো, $\Gamma$
পরিতৃপ্তিযোগ্য হয় ঠিক তখনই, যখন প্রত্যেক সসীম
$\Gamma_0 \subseteq \Gamma$ পরিতৃপ্তিযোগ্য। নির্দিষ্ট ধর্মবিশিষ্ট
গঠনের অস্তিত্ব প্রমাণে সংহতি উপপাদ্য কাজে লাগে। যেমন, এর সাহায্যে
দেখানো যায় যে ইচ্ছামতো বড় সসীম মডেল আছে এমন প্রতিটি তত্ত্বেরই
অসীম মডেল আছে। বিশেষত, এর অর্থ হলো প্রথম-ক্রমের যুক্তিবিদ্যায়
সসীমতা প্রকাশ করা যায় না। দ্বিতীয় অনুসিদ্ধান্ত,
\emph{লোয়েনহাইম--স্কোলেম উপপাদ্য}, বলে যে প্রত্যেক পরিতৃপ্তিযোগ্য
$\Gamma$-র একটি গণনীয় মডেল আছে। এর থেকে আবার দেখা যায় যে
প্রথম-ক্রমের যুক্তিবিদ্যায় অগণনীয়তা প্রকাশ করা যায় না।
